\scanpage{601}
\mspage{604}

Let us examine this central extension
\[
(84)\qquad
\diagram{\begin{tikzcd}[ampersand replacement=\&,column sep=large,row sep=small]
1\arrow[r]\&\mu\arrow[r]\&
\mathcal E\arrow[r,"{(\varepsilon_3,\varepsilon_2)}"]
\arrow[d,phantom,"{\rotatebox{90}{$\simeq$}}" description]\&
\mu\times\mu\arrow[r]\&1\\
\&\&S_0^\Pi\Gamma_{\rho,\sigma}/S_0^{\Pi_0}\Gamma^0_{\rho,\sigma}\&\&
\end{tikzcd}}
\]
fitting into the diagram of extensions,
\[
(85)\qquad
\diagram{\begin{tikzcd}[ampersand replacement=\&,column sep=large,row sep=large]
1\arrow[r,dash]\&\mu\arrow[r]\&
\mathcal E\arrow[r,dash]\&\mu\times\mu\arrow[r,dash]\&1\\
1\arrow[r,dash]\&\mu\arrow[r]\arrow[u,"1"']\arrow[d,hook]\&
S_0^\Pi\Gamma^!_{\rho,\sigma}\arrow[r]\arrow[u]\arrow[d,hook]\&
Z^!_{\rho,\sigma}\arrow[u,"{(\varepsilon_3,\varepsilon_2)}"']
\arrow[d,leftrightarrow,"1"]\arrow[r,phantom,"{\simeq M[0]}" description]\&\vphantom{1}\\
1\arrow[r]\&L_\rho\times L_\sigma\arrow[r]\&
S_0\Gamma^!_{\rho,\sigma}\arrow[r]\&Z^!_{\rho,\sigma}\arrow[r]\&1
\end{tikzcd}}
\]

The interest of (84) is that it allows one to recover the extension
$S_0\Gamma_{\rho,\sigma}$ of $Z^!_{\rho,\sigma}$ by
$L_\rho\times L_\sigma$, by the \unsure{covariant}-\unsure{contravariant}
\unsure{law} in the functor
``extensions'', from
\[
(86)\qquad
\diagram{\begin{tikzcd}[ampersand replacement=\&,column sep=large,row sep=small]
Z^!_{\rho,\sigma}\arrow[r,"{(\varepsilon_3,\varepsilon_2)}"]\&
\mu\times\mu
\end{tikzcd}},\qquad
\begin{gathered}
\diagram{\begin{tikzcd}[ampersand replacement=\&,column sep=large,row sep=small]
\mu\arrow[r,hook]\&L_\rho\times L_\sigma\\
-1\arrow[r,mapsto]\&(\rho^3,\sigma^2)
\end{tikzcd}}\\
\text{(``diagonal homomorphism'')}
\end{gathered},
\]
as well as the subextension $S_0^\Pi\Gamma^!_{\rho,\sigma}$ of
the latter (of $Z^!_{\rho,\sigma}$ by $\mu$).

Let us first study the order of the elements of $\mathcal E$ --- they are
obviously divisors of $4$, are they all of order $2$ (or $1$)? Let
\[
\underline{\unsure{\bar u|u}}=
\left(
\underset{\inserted{\substack{\rotatebox{90}{$\in$}\\Z}}}{u},
\underset{\inserted{\substack{\rotatebox{90}{$\in$}\\N_\rho^*}}}{a},
\underset{\inserted{\substack{\rotatebox{90}{$\in$}\\N_\sigma^*}}}{b}
\right)\in S_0^\Pi\Gamma^!
\qquad\text{i.e.}\quad
\left\{\begin{aligned}
u&\equiv b\quad(\widehat\Pi_0)\\
u&\equiv\varepsilon_3[\sigma]a\quad(\widehat\Pi)
\end{aligned}\right.
\qquad(87)
\]
\margin{\uline{N.B.} To \unsure{prevent|introduce} confusions,
\unsure{I write} $\varepsilon_3[\sigma],\varepsilon_3[\tau']$
in preference \unsure{to|of} $\varepsilon_3(\sigma),\varepsilon_3(\tau')$
for $\sigma^{\unsure{\nu_3}},\tau'^{\unsure{\nu_3}}$.}

I \deleted{\unsure{want}} want to examine whether
$\underline{u}^{2}\in S_0^{\Pi_0}\Gamma^{!0}$.
Obviously
\[
\underline{u}=(u^2,a^2,b^2)\in S_0^\Pi\Gamma^{!0}
\]

\mspage{600}

\scanpage{602}

it remains to see whether $\underline u^{2}\in S_0^{\Pi_0}\Gamma^{!}$, i.e. whether we have
\deleted{\unsure{not}} only
$\unsure{u^{2}}\equiv b^{2}\ (\Pi_0)$ and
$\deleted{\unsure{a^{2}}}\unsure{u^{2}}\equiv a^{2}\ (\widehat\Pi)$
--- which is clear and simply expresses
$\underline u^{2}\in S_0^{\Pi}\unsure{\Gamma|\Gamma^{!}}$ --- but even
\[
\tag{88}
\unsure{u^{2}}\equiv a^{2}\quad\bmod\ \uuline{\widehat\Pi_0}
\qquad (?)
\]

The relation
\[
u\equiv\varepsilon_3[\sigma].a\quad\bmod\Pi
\]
can also be written
\[
u\equiv\pm\varepsilon_3[\sigma].a\quad\bmod\Pi_0
\]
and implies
\[
u^2\equiv\varepsilon_3[\sigma]a\ \varepsilon_3[\sigma]a
\quad\bmod\Pi_0
\]
which we can write
\[
\tag{89}
\underset{\bmod\widehat\Pi_0}{u^2\equiv
\varepsilon_3\,\varepsilon_3[\sigma](a).a}
=\begin{cases}
a^2&\text{if }\varepsilon_3=1.\\
-\sigma(a).a&\text{if }\varepsilon_3=-1
\end{cases}
\]

Thus (88) is true if $\varepsilon_3=1$, whereas
for $\varepsilon_3=-1$, it is equivalent to
\[
\sigma(a)\equiv-a\qquad(\widehat\Pi_0)
\]
i.e.
\[
[\sigma,a]^{\cdot}=-1\quad\text{in }\unsure{\widetilde{\mathcal G}|\widetilde{\mathcal C}}/\Pi_0
\]
i.e.
\[
\underbrace{\varepsilon_2(a)}_{\inserted{\begin{gathered}\rotatebox{90}{$=$}\\\varepsilon_2(u)\end{gathered}}}=-1
\qquad\text{i.e. }\varepsilon_2(\unsure{u})\varepsilon_3(\unsure{u})=1.
\]

Thus we find:

\uline{Proposition} Let $x\in\mathcal E$ (cf. (84)). For
$x^2=1$ to hold, it is necessary and sufficient that we have
$\varepsilon_3(x)=1$ \uline{or} $\varepsilon_2(x)\varepsilon_3(x)=1$,
i.e. that
\deleted{$x$ \unsure{belong to the subgroup} $\mathcal E'$ \unsure{or}}
\inserted{we have}
\[
\tag{90}
x\in\mathcal E'\cup\mathcal E''',
\]

\mspage{601}

\scanpage{603}
\mspage{605}

where
\[
\tag{91}
\mathcal E'=\{\unsure{x|n}\in\mathcal E\mid\varepsilon_3(x)=1\},\qquad
\mathcal E'''=\{\unsure{x|n}\in\mathcal E\mid\varepsilon_2(x)\varepsilon_3(x)=1\}
\]

\uline{Corollary 1} $\mathcal E$ has exactly two elements
of order $4$, which are the two elements of $\mathcal E$
above the element $(-1,1)$ of $\mu\times\mu$.
These two elements are inverses of one
another, \unsure{these|and} are the two generators of
the only cyclic subgroup of order $4$ of $\mathcal E$,
which is none other than
\[
\tag{92}
\mathcal E''=\{\unsure{x|n}\in\mathcal E\mid\varepsilon_2(x)=1\}
\]

\uline{Corollary 2} The extension $\mathcal E$ is not split.

For being central, it \deleted{\unsure{itself}} would be isomorphic
to $(\mu\times\mu)\times\mu$, so its elements would have order
dividing $2$.

We can now regard
$\mathcal E$ as an extension of $\mathcal E/\mathcal E''\simeq\{\pm1\}=\mu$ by
$\mathcal E''\ (\simeq\mathbb Z/4\mathbb Z)$.
\deleted{But it} It would be necessary
first to make explicit the operation of $\mathcal E/\mathcal E''=\mu$
on $\unsure{\mathcal E'|\mathcal E''}$ --- the element \inserted{$-1$} of order $2$ of $\mu$
induces an involutive \unsure{automorphism} of \deleted{$\unsure{\mathbb Z/4\mathbb Z}$} $\mathcal E''$
\inserted{moreover in $\mathrm{Aut}(\unsure{\mathcal E'|\mathcal E''})\simeq
(\mathbb Z/4\mathbb Z)^*\simeq\{\pm1\}$, so
the automorphism in question}
\deleted{that \unsure{its}} is multiplication by a
sign $\varepsilon\ \pm1$, which has to be determined. We are going
to take \inserted{as representative of $-1\in\mathcal E/\mathcal E''$}
the element of $\mathcal E$ coming
\deleted{from}\inserted{from}
\[
\tag{92}
\deleted{\unsure{\tau,a',\tau}}\quad
\underline x=(u_0=\tau,\ a_0=\tau',\ b_0=\tau)
\in S_0^{\Pi_0}\Gamma^{!'''}\qquad(\text{cf.}(70))
\]
which we are going to have operate on the image of a
\[
\tag{93}
\underline u=(u,a,b)\in S_0^{\Pi}\Gamma^{!},\quad
\deleted{\unsure{\text{where}}}\ \text{with }\varepsilon_3(u)=-1,\ \varepsilon_2(u)=1
\]
\inserted{\[
\left(\begin{array}{ll}
u\equiv b\ \deleted{\unsure{\equiv a}}&\bmod\widehat\Pi_0\\
u\equiv\sigma a&\bmod\widehat\Pi
\end{array}\right)
\]}
(which is a generator of $\mathcal E''$).

\mspage{602}

\scanpage{604}

We shall thus have
\[
\underline x(\underline u)=(\tau(u),\tau'(a),\tau(b))\equiv
\deleted{\unsure{(u,\varepsilon a,\varepsilon b)|(\varepsilon u,\varepsilon a,\varepsilon b)}}
\underbrace{\inserted{(1,\varepsilon,\varepsilon)(u,a,b)}}_{\inserted{(u,\varepsilon a,\varepsilon b)}}
\]
modulo $S_0^{\Pi_0}\Gamma^{\unsure{!0|0}}$, with a sign
$\varepsilon\in\{\pm1\}$ to be determined. The congruence condition can be written
\[
(u\tau(u)^{-1},\ \varepsilon a\tau'(a)^{-1},\
\varepsilon b\unsure{\tau'(b)|\tau(b)}^{-1})\in S_0^{\Pi_0}\Gamma^{\unsure{!0|0}}
\]
that we are in $S_0^\Pi\Gamma^{!0}$ is obvious,
\deleted{\unsure{the choice}}
\deleted{\unsure{immaterial|independent}}
\inserted{\unsure{valid|valid}}
whatever the choice of the sign
$\varepsilon$, \deleted{\unsure{it is}}
\deleted{\unsure{it is|and}}\inserted{\unsure{it is}}
\unsure{the|the} \unsure{condition|conditions} of congruence
\deleted{which} \deleted{\unsure{are essential}}
\inserted{which \unsure{distinguishes|distinguish} $S_0^{\Pi_0}\Gamma^{!}$ from
$S_0^\Pi\Gamma^{!}$ which is}
essential, namely
\[
\tag{94}
u\tau(u)^{-1}\equiv\varepsilon\unsure{a}\tau'(\unsure{a})^{-1}
\qquad(\widehat\Pi_0)\qquad\text{i.e.}
\]
\deleted{\[
\unsure{u(\tau)\tau^{-1}\equiv\varepsilon a(\tau')\tau'^{-1}}
\]}

The difficulty with this approach comes
from the fact that the only elements of
$M$ that we know how to make explicit are those
coming from $\unsure{\widetilde{\mathcal G}|\widetilde{\mathcal C}}\simeq\mathrm{Gl}(2,\mathbb Z)$
itself; now
for these we have $\varepsilon_2(u)\varepsilon_3(u)=1$, whereas
the $u$-\inserted{here} \deleted{\unsure{test}} above must satisfy
$\varepsilon_2(u)=1$, $\varepsilon_3(u)=-1$ so
$\varepsilon_2(u)\varepsilon_3(u)=-1$.
We \unsure{can} on the other hand work in the
\deleted{\unsure{\text{group}} $\unsure{\widetilde{\mathcal G}|\widetilde{\mathcal C}}$}
\inserted{quotient} $\mathrm{Gl}(2,\widehat{\mathbb Z})$ of
\deleted{$\unsure{\widetilde{\mathcal G}|\widetilde{\mathcal C}}$} $M$, and
take in it
\[
\tag{95}
u=\begin{pmatrix}\mu&0\\0&1\end{pmatrix}
\quad\text{with}\quad
\mu\in\widehat{\mathbb Z}^{*},\quad
\begin{cases}
\mu\equiv-1&(3)\\
\mu\equiv1&(4)
\end{cases}
\]
\inserted{\[\text{i.e.}\quad\mu\equiv5\quad(12)\]}
whereas $\tau$ corresponds to the
matrix
\[
\tau=\begin{pmatrix}-1&0\\0&1\end{pmatrix}
\]

\margin{$\unsure{\wedge}$}
\mspage{603}

\scanpage{605}
\mspage{606}

we shall in any case have
\[
u\tau(u)^{-1}=1
\]
let, on the other hand, in $\mathrm{Gl}(2,\widehat{\mathbb Z})$
\[
\tag{96}
\left\{\begin{array}{lll}
a=\tau'a_0&a_0=c\rho+d&c^2+cd+d^2\in\widehat{\mathbb Z}^{*}\\[3pt]
&\tau'=\begin{pmatrix}0&1\\1&0\end{pmatrix},&
\rho=\begin{pmatrix}0&1\\-1&1\end{pmatrix}
\end{array}\right.
\]
\unsure{thus|whence}
\[
\tag{97}
a=\begin{pmatrix}0&1\\1&0\end{pmatrix}
\begin{pmatrix}d&c\\-c&c+d\end{pmatrix}
=\begin{pmatrix}-c&c+d\\d&c\end{pmatrix}
\]
and
\[
\tag{98}
\mu=-(c^2+cd+d^2)
\]

\margin{\unsure{a little brief}... cf. below...}

We \deleted{\unsure{are going to}} work with the closed \inserted{sub-}group
$\widehat\Pi_0$ in $\mathrm{Sl}(2,\widehat{\mathbb Z})$ generated by
\[
\begin{gathered}
-\ell_0\underset{\unsure{,}}{\unsure{=}}\unsure{-}\begin{pmatrix}
1&-2\\\deleted{\unsure{-2|2}}0&1
\end{pmatrix},\qquad
-\ell_1=-\begin{pmatrix}1&0\\2&1\end{pmatrix},\qquad
-\ell_\infty=-\begin{pmatrix}3&-2\\2&-1\end{pmatrix}
\end{gathered}
\]
and in the quotient group of
$\mathrm{Gl}(2,\widehat{\mathbb Z})$ by $\widehat\Pi_0$, which is
\deleted{\unsure{the}} the extension of
$\widehat{\mathbb Z}^{*}$ by
\[
\mathrm{Sl}(2,\widehat{\mathbb Z})/\widehat\Pi_0
\simeq\mathrm{Sl}(2,\mathbb Z)/\Pi_0
\simeq\left\{(\rho,\sigma)\ \left|\begin{array}{l}
\rho^6=\sigma^4=1\\
\rho^3=\sigma^2\\
\sigma(\rho)=\rho^{-1}
\end{array}\right.\right\},
\]
\deleted{\unsure{with}} and \unsure{this|in} a $\frac12$ direct product thanks to
\[
T_0=\left\{\begin{pmatrix}
\deleted{\unsure{\xi}}\unsure{\zeta|\xi}&0\\0&1
\end{pmatrix}\ \middle|\ \unsure{\zeta|\xi}\in\widehat{\mathbb Z}^{*}\right\}
\overset{\det}{\simeq}\widehat{\mathbb Z}^{*},
\]
operating through $\varepsilon_2(\unsure{\zeta|\xi})$
--- cf. pages $593'$, $594$ --- formulas (34)(35)(35 bis) ].
We \inserted{must} suppose
\[
\tag{99}
\sigma a\equiv u\qquad(\widehat\Pi_{\unsure{0}})
\qquad\text{where }\sigma=\begin{pmatrix}0&-1\\1&0\end{pmatrix}
\]
thus
\[
\tag{100}
\sigma a=\begin{pmatrix}-d&-c\\-c&c+d\end{pmatrix}
\]

\mspage{604}